On considère un solide en rotation autour d'un axe fixe, avec une fréquence
$f=\qty{0.04}{\hertz}$.

Soit un point $M$ du solide qui décrit une trajectoire circulaire de rayon
$r=\qty{2}{\cm}$ et de centre qui appartient a l'axe de rotation.

À l'instant $t=\qty{1}{\second}$, le point $M$ est repéré par l'angle
$\theta=\ang{30}$.

\begin{enumerate}
    \item Calculer la période du mouvement.
    \item Calculer la valeur de la vitesse angulaire du mouvement, et la valeur
          de la vitesse linéaire du point $M$.
    \item Écrire l'équation horaire du mouvement.
    \item Déterminer le nombre de tours qu'a effectué le point $M$ a l'instant
          $t^{\prime}=\qty{25,135}{\second}$.
    \item À quelle date le point $M$ fait son deuxième passage par l'origine des
          abscisses angulaire.
\end{enumerate}
